\documentclass[11pt]{article}
\usepackage[margin=1in]{geometry}
\usepackage{amsmath,amssymb,amsthm}
\usepackage{xurl}
\usepackage{hyperref}
\sloppy
\let\oldus\_ \renewcommand{\_}{\oldus\allowbreak}
\newtheorem{theorem}{Theorem}
\newtheorem{lemma}[theorem]{Lemma}
\theoremstyle{remark}
\newtheorem{remark}{Remark}

\title{Conjecture 00000004011 is false for linear-growth vector fields:\\
the linear RDE $dY = Y\,dX$ with a Gaussian driver has lognormal, not Gaussian, tails}
\author{Jackmeson1}
\date{}

\begin{document}
\maketitle

\section{The conjecture and our reading}

\textbf{Statement.} \emph{Definition: An RDE is a rough differential equation $dY = V(Y)\,dX$.
Conjecture: The tail of the solution norm satisfies the Fernique-type bound
$P(\|Y\| > t) \le 2\exp(-t^2/c)$, where the constant $c$ depends only on the growth constant of
the vector fields and the tail parameters of the driver.}

\textbf{Reading.} We read ``growth constant of the vector fields'' as the constant $K$ in the
linear growth condition $|V(y)| \le K(1+|y|)$. Linear-growth RDEs are a standard class in the
integrability literature: Friz and Riedel (arXiv:1104.0577) write in their abstract
``Integrability properties of (classical, linear, linear growth) rough differential equations
(RDEs) are considered'', and their paper has a ``Remark 5 (Unbounded vector fields)''.
The claim is: there is $c > 0$, determined by $K$ and by the tail parameters of the driver, such
that $P(\|Y\| > t) \le 2\exp(-t^2/c)$ for all $t \ge 0$. We refute it for one equation with $K = 1$
and one Gaussian driver, so no choice of the function $(K,\text{tail parameters}) \mapsto c$ can
work. We also refute the weaker version that asks for the bound only for $t \ge T$, for any $T$.

\textbf{Reading not refuted.} If ``growth constant'' is read as a uniform bound
$\sup_y |V(y)|$ (bounded vector fields), then $V(y) = y$ is out of scope and this note says
nothing. In that setting Gaussian tails are expected: Friz and Riedel recall that
``we recover the well-known fact that solutions of the Stratonovich SDE have Gaussian tails at any
fixed point provided'' the vector fields are sufficiently smooth.

\textbf{The counterexample.}
\begin{itemize}
\item Vector field $V(y) = y$ on $\mathbb{R}$. It satisfies $|V(y)| \le 1\cdot(1+|y|)$ (growth
constant $1$) and is $1$-Lipschitz.
\item Driver $X_t = tG$ on $[0,1]$, where $G$ is a standard Gaussian random variable on a
probability space $(\Omega, P)$, i.e.\ $P \circ G^{-1} = N(0,1)$. Each sample path is smooth, and
$|X_t - X_s| = |G|\,|t-s|$, so the sup norm, the $p$-variation norm and the Hoelder seminorms of $X$
on $[0,1]$ all equal $|G|$. The driver has Gaussian tails: $P(|G| > t) \le 2\exp(-t^2/2)$ for
$t \ge 0$. So its tail parameters are those of a standard Gaussian, and the driver itself satisfies
the Fernique-type bound with $c = 2$.
\item Initial value $Y_0 = 1$.
\end{itemize}

\textbf{RDEs with smooth drivers.} Rough path theory extends classical controlled differential
equations: Wikipedia (``Rough path'') says ``In stochastic analysis, a rough path is a
generalization of the classical notion of a smooth path'' and ``This makes it possible to define
and solve controlled differential equations of the form'' $dy_t = f(y_t)\,dx_t$. In its version of
``Lyons' Universal Limit theorem'', $X(n)$ is ``a sequence of paths with finite total variation''
and $Y(n)$ is ``the solution to the differential equation'' $dY(n) = V(Y(n))\,dX(n)$ in the
classical sense; the rough-path solution is the limit of these classical solutions. For a driver
that is $C^1$ on $[0,1]$ we therefore use the classical equation
\[
  Y(0) = y_0, \qquad Y'(t) = V(Y(t))\,X'(t) \quad (t \in [0,1]),
\]
with one-sided derivatives at $0$ and $1$ (Lean: \texttt{HasDerivWithinAt} on $[0,1]$, with
$X'$ given by \texttt{derivWithin X (Icc 0 1)}). This is the definition \texttt{IsRDESolution}.

\textbf{Conventions.} For $0<\alpha\le1$ the $\alpha$-Hoelder \emph{seminorm} of the driver equals
$|G|$; the full Hoelder norm (sup norm plus seminorm) equals $2|G|$, which changes only a fixed
constant in the driver's Gaussian tail. For an arbitrary functional $N$ whose tail events need not be
measurable, the general Lean conclusion is an outer-measure statement; for the concrete sup norm,
$N(Y)=\max(1,e^{G})$ is measurable and the statement is an ordinary probability bound.

\textbf{The solution norm.} The text does not fix $\|Y\|$. We allow any path functional $N$ with
$Y_1 - Y_0 \le N(Y)$ on the solution paths. This includes the sup norm on $[0,1]$ (proved in Lean),
the terminal value $|Y_1|$ (because $Y_0 = 1 > 0$), and the $p$-variation and Hoelder norms on
$[0,1]$, each of which is at least $|Y_1 - Y_0|$.

\textbf{Conventions.} $N(0,1)$ is Mathlib's \texttt{gaussianReal 0 1}, with density
$\varphi(x) = (2\pi)^{-1/2}e^{-x^2/2}$. Probabilities are $[0,\infty]$-valued measures and
$2\exp(-t^2/c)$ enters as \texttt{ENNReal.ofReal}. The tail event is the strict one
$\{N(Y) > t\}$. The probability space is arbitrary; we only assume that $G$ is measurable with law
$N(0,1)$. Measurability of $\omega \mapsto N(Y_\omega)$ is not assumed (the lower bound uses only
monotonicity of the measure).

\section{Proof}

\begin{lemma}[existence and uniqueness]
For every $g \in \mathbb{R}$, $t \mapsto e^{tg}$ solves $Y' = Y X'$, $Y_0 = 1$ with
$X_t = tg$ on $[0,1]$. Every solution $Y$ satisfies $Y(t) = e^{tg}$ for $t \in [0,1]$.
\end{lemma}
\begin{proof}
Existence: $X' = g$ on $[0,1]$ and $(e^{tg})' = g e^{tg}$. Uniqueness: $f(t) = Y(t)e^{-tg}$ has
$f' = gYe^{-tg} - gYe^{-tg} = 0$ on $[0,1]$, so $f$ is constant, $f(t) = f(0) = 1$.
\end{proof}

So for every $\omega$, $Y_\omega(1) - Y_\omega(0) = e^{G(\omega)} - 1$, a shifted lognormal variable ($e^{G}$ itself is lognormal).

\begin{lemma}[Gaussian tail lower bound]
For $s \ge 0$, $P(G > s) \ge \varphi(s+1)$.
\end{lemma}
\begin{proof}
$P(G > s) \ge \int_{s}^{s+1}\varphi(x)\,dx \ge \varphi(s+1)$, since $\varphi$ is decreasing on
$[0,\infty)$.
\end{proof}

\begin{theorem}[main]
Let $Y_\omega$ be any solution of $dY = Y\,dX$, $Y_0 = 1$ with driver $X_t = tG(\omega)$, and let
$N$ satisfy $Y_\omega(1) - Y_\omega(0) \le N(Y_\omega)$ for all $\omega$. For every $c > 0$ and every
$T \in \mathbb{R}$ there is $t \ge T$ with $P(N(Y) > t) > 2\exp(-t^2/c)$.
\end{theorem}
\begin{proof}
Put $u = 1 + 20c + \max(T,0)$ and $t = e^u - 1$. Then $u \ge 1$, $t \ge u \ge T$, and
$t \ge u^2/2$ because $e^u \ge 1 + u + u^2/2$. If $G(\omega) > u$ then
$N(Y_\omega) \ge e^{G(\omega)} - 1 > e^u - 1 = t$. Hence
$P(N(Y) > t) \ge P(G > u) \ge \varphi(u+1)$. It remains to show
$2e^{-t^2/c} < (2\pi)^{-1/2}e^{-(u+1)^2/2}$. We have $u^2 > 20c$, $t^2 \ge u^4/4$ and
$(u+1)^2/2 + 3 \le 5u^2$, so
$((u+1)^2/2 + 3)\,c \le 5u^2 c < u^4/4 \le t^2$, i.e.\ $3 - t^2/c < -(u+1)^2/2$. Since
$\sqrt{2\pi} < 3$ and $e^3 \ge 1 + 3 + 9/2 > 6$,
$2\sqrt{2\pi}\,e^{-t^2/c} < 6e^{-t^2/c} < e^{3 - t^2/c} < e^{-(u+1)^2/2}$.
\end{proof}

\begin{remark}
Taking $T = 0$: no $c > 0$ gives $P(N(Y) > t) \le 2\exp(-t^2/c)$ for all $t \ge 0$. This is
consistent with the literature: for linear RDEs, Friz and Riedel (arXiv:1104.0577, Remark 7) note,
in the case ``which covers Brownian driving signals, we have log-normal tails. This is in agreement
with trivial examples such as the standard Black-Scholes model in which the stock price is
log-normally distributed.'' The same mechanism works for the Brownian driver
($dY = Y\circ dB$, $Y_t = e^{B_t}$); we formalize only the smooth Gaussian driver.
\end{remark}

\section{Lean formalization}

File \texttt{lean/Conjecture4011/Basic.lean} (namespace \texttt{C4011}, only
\texttt{import Mathlib}, 281 lines). Definitions: \texttt{HasLinearGrowth V K}
($\forall y,\ |V y| \le K(1+|y|)$), \texttt{IsRDESolution V X y0 Y}, \texttt{linDriver g}
($t \mapsto t g$), \texttt{supNorm01 y} ($\sup_{t\in[0,1]}|y(t)|$ as an \texttt{sSup}). Theorems:
\begin{itemize}
\item \texttt{linearGrowth\_id : HasLinearGrowth id 1}; \texttt{lipschitz\_id\_field};
\texttt{linDriver\_increment}: $|X_t - X_s| = |g||t-s|$.
\item \texttt{exp\_isRDESolution}, \texttt{isRDESolution\_unique} (Lemma 1).
\item \texttt{gaussian\_tail\_lower} (Lemma 2); \texttt{driver\_tail}:
$P\{t < |G|\} \le 2e^{-t^2/2}$ for $t \ge 0$ (Chernoff bound via Mathlib's
\texttt{HasSubgaussianMGF.measure\_ge\_le}).
\item \texttt{tail\_not\_fernique} (Theorem 3): for $\Omega$, $P$, measurable $G$ with
\texttt{P.map G = gaussianReal 0 1}, any $Y$ with
\texttt{$\forall\omega$, IsRDESolution id (linDriver (G $\omega$)) 1 (Y $\omega$)} and any $N$ with
\texttt{$\forall\omega$, Y $\omega$ 1 - Y $\omega$ 0 $\le$ N (Y $\omega$)}:
\texttt{$\forall$ c, 0 < c $\to$ $\forall$ T, $\exists$ t $\ge$ T, ENNReal.ofReal (2 * exp (-t\^{}2/c)) < P \{$\omega$ | t < N (Y $\omega$)\}}.
\item \texttt{no\_fernique\_constant}: under the same hypotheses,
\texttt{$\neg\,\exists$ c, 0 < c $\wedge$ $\forall$ t, 0 $\le$ t $\to$ P \{$\omega$ | t < N (Y $\omega$)\} $\le$ ENNReal.ofReal (2 * exp (-t\^{}2/c))}.
\item \texttt{no\_fernique\_constant\_supNorm}: the same with $N = $ \texttt{supNorm01}.
\item \texttt{concrete\_instance}: on $(\mathbb{R}, N(0,1))$ with $G = \mathrm{id}$ and
$Y_\omega(t) = e^{t\omega}$ all hypotheses hold and the sup-norm bound fails.
\end{itemize}
\textbf{Axiom audit.} \texttt{lake build} succeeds. \texttt{\#print axioms} for
\texttt{tail\_not\_fernique}, \texttt{no\_fernique\_constant}, \texttt{no\_fernique\_constant\_supNorm},
\texttt{concrete\_instance}, \texttt{driver\_tail} and \texttt{isRDESolution\_unique} gives
\texttt{[propext, Classical.choice, Quot.sound]}. There is no \texttt{sorry}, no
\texttt{native\_decide} and no added axiom.

\section*{Sources}
\begin{itemize}
\item Wikipedia, ``Rough path'', \url{https://en.wikipedia.org/wiki/Rough_path} (introduction and
section ``Universal limit theorem'').
\item P. Friz, S. Riedel, ``Integrability of (non-)linear rough differential equations and
integrals'', arXiv:1104.0577, \url{https://arxiv.org/abs/1104.0577} (abstract; Remarks 5, 6, 7).
\end{itemize}

\end{document}
